\documentclass[12pt,a4paper]{article}

\usepackage{amsmath}
\usepackage{amssymb}
\usepackage{mathtools}

\title{Algebraic Representation of the Zebra Program}
\author{}
\date{}

\begin{document}
\maketitle

\section*{Algebraic Operators}

We assume the following algebraic operators:

\begin{itemize}
  \item $\circ$ : sequential composition of constraints
  \item $\otimes$ : parallel (independent) composition
  \item $\oplus$ : nondeterministic choice
  \item $\mu$ : least fixpoint (recursion / search)
\end{itemize}

Primitive constraint generators:
\[
\mathrm{unify}(\cdot),\quad
\mathrm{mem}(\cdot),\quad
\mathrm{shift},\quad
\mathrm{swap}
\]

\section*{Definition of \texttt{iright}}

The right-adjacency relation is defined as a recursive constraint generator:

\[
\mathrm{ir}
\;:=\;
\mu X.\;
\Big(
  \mathrm{match}(L,R)
  \;\oplus\;
  \mathrm{shift} \circ X
\Big)
\]

\section*{Definition of \texttt{nextto}}

The symmetric adjacency relation is defined by disjunction:

\[
\mathrm{nt}
\;:=\;
\mathrm{ir}
\;\oplus\;
\mathrm{swap} \circ \mathrm{ir}
\]

\section*{Algebraic Form of the Zebra Program}

The Zebra puzzle is represented as a single fixpoint expression:

\[
\begin{aligned}
Z
\;:=\;
\mu\Bigg(
&
\mathrm{unify}\Big(
[\,\mathrm{house}(\mathrm{norwegian},\_,\_,\_,\_),
\_,
\mathrm{house}(\_,\_,\_,\mathrm{milk},\_),
\_,
\_]
\Big)
\\[4pt]
&\circ
\bigotimes
\begin{aligned}[t]
(&\mathrm{mem}(\mathrm{house}(\mathrm{englishman},\_,\_,\_,\mathrm{red})) \\
\otimes\;&\mathrm{mem}(\mathrm{house}(\mathrm{spaniard},\mathrm{dog},\_,\_,\_)) \\
\otimes\;&\mathrm{mem}(\mathrm{house}(\_,\_,\_,\mathrm{coffee},\mathrm{green})) \\
\otimes\;&\mathrm{mem}(\mathrm{house}(\mathrm{ukrainian},\_,\_,\mathrm{tea},\_)) \\
\otimes\;&\mathrm{mem}(\mathrm{house}(\_,\mathrm{snails},\mathrm{winston},\_,\_)) \\
\otimes\;&\mathrm{mem}(\mathrm{house}(\_,\_,\mathrm{kools},\_,\mathrm{yellow})) \\
\otimes\;&\mathrm{mem}(\mathrm{house}(\_,\_,\mathrm{luckystrike},\mathrm{oj},\_)) \\
\otimes\;&\mathrm{mem}(\mathrm{house}(\mathrm{japanese},\_,\mathrm{parliaments},\_,\_)) \\
\otimes\;&\mathrm{mem}(\mathrm{house}(W,\_,\_,\mathrm{water},\_)) \\
\otimes\;&\mathrm{mem}(\mathrm{house}(Z,\mathrm{zebra},\_,\_,\_))
)
\end{aligned}
\\[6pt]
&\circ
\Big(
\mathrm{ir}\big(
\mathrm{house}(\_,\_,\_,\_,\mathrm{ivory}),
\mathrm{house}(\_,\_,\_,\_,\mathrm{green})
\big)
\\
&\qquad\otimes\;
\mathrm{nt}\big(
\mathrm{house}(\_,\_,\mathrm{chesterfield},\_,\_),
\mathrm{house}(\_,\mathrm{fox},\_,\_,\_)
\big)
\\
&\qquad\otimes\;
\mathrm{nt}\big(
\mathrm{house}(\_,\_,\mathrm{kools},\_,\_),
\mathrm{house}(\_,\mathrm{horse},\_,\_,\_)
\big)
\\
&\qquad\otimes\;
\mathrm{nt}\big(
\mathrm{house}(\mathrm{norwegian},\_,\_,\_,\_),
\mathrm{house}(\_,\_,\_,\_,\mathrm{blue})
\big)
\Big)
\Bigg)
\end{aligned}
\]

\section*{Interpretation}

This single expression represents the complete Zebra program.
Execution corresponds to projection into a lower-dimensional operational
semantics (e.g.\ Prolog), while refactoring and optimization correspond
to algebraic transformation at this level.

\end{document}
